% Source: https://artofproblemsolving.com/community/c3h1868037_intermediate_counting_and_probability_problem_1135
% Stable identity: aops-topic-1868037-q01
A black ball can only come from urn B, so the second draw must be from urn B (probability $\dfrac{1}{2}$). What urn B contains at that moment depends on the first draw.
\begin{itemize}
\item First draw from A (probability $\dfrac{1}{2}$): urn B still has $3$ red and $3$ black balls, so $P(\text{2nd black})=\dfrac{1}{2}\cdot\dfrac{3}{6}=\dfrac{1}{4}$. The first ball is red with probability $\dfrac{2}{6}=\dfrac{1}{3}$.
\item First draw from B and red (probability $\dfrac{1}{2}\cdot\dfrac{3}{6}=\dfrac{1}{4}$): urn B has $2$ red and $3$ black, so $P(\text{2nd black})=\dfrac{1}{2}\cdot\dfrac{3}{5}=\dfrac{3}{10}$.
\item First draw from B and black (probability $\dfrac{1}{4}$): urn B has $3$ red and $2$ black, so $P(\text{2nd black})=\dfrac{1}{2}\cdot\dfrac{2}{5}=\dfrac{1}{5}$.
\end{itemize}
Hence
\[
P(\text{1st red and 2nd black})=\dfrac{1}{2}\cdot\dfrac{1}{3}\cdot\dfrac{1}{4}+\dfrac{1}{4}\cdot\dfrac{3}{10}=\dfrac{1}{24}+\dfrac{3}{40}=\dfrac{14}{120}=\dfrac{7}{60},
\]
\[
P(\text{2nd black})=\dfrac{1}{2}\cdot\dfrac{1}{4}+\dfrac{1}{4}\cdot\dfrac{3}{10}+\dfrac{1}{4}\cdot\dfrac{1}{5}=\dfrac{1}{8}+\dfrac{3}{40}+\dfrac{1}{20}=\dfrac{15+9+6}{120}=\dfrac{1}{4} .
\]
By the definition of conditional probability,
\[
P(\text{1st red}\mid\text{2nd black})=\dfrac{\dfrac{7}{60}}{\dfrac{1}{4}}=\boxed{\dfrac{7}{15}}.
\]
